% GENERATED FILE. Edit canonical lessons, then run npm run generate:books.
% canonical-source-hash: fnv1a64:6adfa973c8ce9fba
% canonical-lessons: TE-C403-paripalana, TE-C403-pratinidhi, TE-C403-samanvayam, TE-C403-pathakam, TE-C403-sanksemam

\chapter{Running Things}
\label{ch:te-running-things}
\hlchaptermodality{403}

\begin{chapteropening}
\textbf{By the end of this chapter:} \emph{I can talk about administration, schemes and welfare.}

Complete the last lesson of chapter 403: I can talk about administration, schemes and welfare.
\end{chapteropening}

\section[\texorpdfstring{\te{పరిపాలన} (paripālana)}{పరిపాలన (paripālana)}]{\te{పరిపాలన} (paripālana) — administration; governance}

\label{lesson:TE-C403-paripalana}



\begin{quote}
Before the new one: say the Telugu for to govern, then the Telugu for to approve.
\end{quote}

\subsection*{You'll want to know: \te{పరిపాలన}}
\textbf{\te{పరిపాలన}} — \emph{paripālana} — \textquotedblleft{}administration; governance\textquotedblright{}.

\begin{quote}
\te{మంచి} \te{పరిపాలన} \te{వల్ల} \te{జనం} \te{సంతోషంగా} \te{ఉన్నారు}.
\end{quote}

\begin{quote}
\emph{mañci paripālana valla janaṁ santōṣaṅgā unnāru.}
\end{quote}

\begin{quote}
Thanks to good administration the people are happy.
\end{quote}

From the same word as \te{పాలించు}, with \te{పరి}-, all round.

\begin{cousinweb}
\textbf{\te{పరిపాలన}} (\emph{paripālana}) is borrowed from Sanskrit \emph{paripālana}, \textquotedblleft{}the act of guarding; fostering\textquotedblright{} (Monier-Williams), \emph{pari-}, \textquotedblleft{}all round\textquotedblright{}, added to \emph{pālana}.
\end{cousinweb}

\begin{grammarlens}[title={What you've built}]
One more word for giving reasons and opinions.
\end{grammarlens}

\subsection*{Your turn}
\begin{itemize}
\raggedright
  \item \emph{Say it:} \emph{paripālana}
  \item \emph{Say it:} \emph{paripālana}, once more
  \item \emph{Say it:} \emph{mañci paripālana valla janaṁ santōṣaṅgā unnāru.} — Thanks to good administration the people are happy.
  \item \emph{Recall:} say the Telugu for to govern, then the Telugu for to approve, then say \emph{paripālana} again
\end{itemize}

\begin{tcolorbox}[breakable,colback=teal!4,colframe=teal!35!black,title={Before you move on}]
What does \te{పరిపాలన} mean? (\textquotedblleft{}Administration; governance\textquotedblright{}.) Say it once more.
\end{tcolorbox}
\section[\texorpdfstring{\te{ప్రతినిధి} (pratinidhi)}{ప్రతినిధి (pratinidhi)}]{\te{ప్రతినిధి} (pratinidhi) — a representative; a delegate}

\label{lesson:TE-C403-pratinidhi}



\begin{quote}
Before the new one: say the Telugu for to approve, then the Telugu for administration.
\end{quote}

\subsection*{You'll want to know: \te{ప్రతినిధి}}
\textbf{\te{ప్రతినిధి}} — \emph{pratinidhi} — \textquotedblleft{}a representative; a delegate\textquotedblright{}.

\begin{quote}
\te{మా} \te{కంపెనీ} \te{ప్రతినిధి} \te{సమావేశానికి} \te{వెళ్ళాడు}.
\end{quote}

\begin{quote}
\emph{mā kampenī pratinidhi samāvēśāniki veḷḷāḍu.}
\end{quote}

\begin{quote}
Our company’s representative went to the meeting.
\end{quote}

Plural: \te{ప్రతినిధులు}, \emph{pratinidhulu}.

Plural \te{ప్రతినిధులు}.

\begin{cousinweb}
\textbf{\te{ప్రతినిధి}} (\emph{pratinidhi}) is borrowed from Sanskrit \emph{pratinidhi}, \textquotedblleft{}a substitute, a representative, a proxy\textquotedblright{} (Monier-Williams).
\end{cousinweb}

\begin{grammarlens}[title={What you've built}]
One more word for giving reasons and opinions.
\end{grammarlens}

\subsection*{Your turn}
\begin{itemize}
\raggedright
  \item \emph{Say it:} \emph{pratinidhi}
  \item \emph{Say it:} \emph{pratinidhi}, once more
  \item \emph{Say it:} \emph{mā kampenī pratinidhi samāvēśāniki veḷḷāḍu.} — Our company’s representative went to the meeting.
  \item \emph{Recall:} say the Telugu for to approve, then the Telugu for administration, then say \emph{pratinidhi} again
\end{itemize}

\begin{tcolorbox}[breakable,colback=teal!4,colframe=teal!35!black,title={Before you move on}]
What does \te{ప్రతినిధి} mean? (\textquotedblleft{}A representative; a delegate\textquotedblright{}.) Say it once more.
\end{tcolorbox}
\section[\texorpdfstring{\te{సమన్వయం} (samanvayaṁ)}{సమన్వయం (samanvayaṁ)}]{\te{సమన్వయం} (samanvayaṁ) — coordination}

\label{lesson:TE-C403-samanvayam}



\begin{quote}
Before the new one: say the Telugu for administration, then the Telugu for a representative.
\end{quote}

\subsection*{You'll want to know: \te{సమన్వయం}}
\textbf{\te{సమన్వయం}} — \emph{samanvayaṁ} — \textquotedblleft{}coordination\textquotedblright{}.

\begin{quote}
\te{రెండు} \te{శాఖల} \te{మధ్య} \te{సమన్వయం} \te{లేదు}.
\end{quote}

\begin{quote}
\emph{reṇḍu śākhala madhya samanvayaṁ lēdu.}
\end{quote}

\begin{quote}
There is no coordination between the two departments.
\end{quote}

\te{సమన్వయం} \te{చేయు}, to coordinate.

\begin{cousinweb}
\textbf{\te{సమన్వయం}} (\emph{samanvayaṁ}) is borrowed from Sanskrit \emph{samanvaya}, \textquotedblleft{}regular succession, connected sequence, conjunction\textquotedblright{} (Monier-Williams).
\end{cousinweb}

\begin{grammarlens}[title={What you've built}]
One more word for giving reasons and opinions.
\end{grammarlens}

\subsection*{Your turn}
\begin{itemize}
\raggedright
  \item \emph{Say it:} \emph{samanvayaṁ}
  \item \emph{Say it:} \emph{samanvayaṁ}, once more
  \item \emph{Say it:} \emph{reṇḍu śākhala madhya samanvayaṁ lēdu.} — There is no coordination between the two departments.
  \item \emph{Recall:} say the Telugu for administration, then the Telugu for a representative, then say \emph{samanvayaṁ} again
\end{itemize}

\begin{tcolorbox}[breakable,colback=teal!4,colframe=teal!35!black,title={Before you move on}]
What does \te{సమన్వయం} mean? (\textquotedblleft{}Coordination\textquotedblright{}.) Say it once more.
\end{tcolorbox}
\section[\texorpdfstring{\te{పథకం} (pathakaṁ)}{పథకం (pathakaṁ)}]{\te{పథకం} (pathakaṁ) — a scheme; a programme}

\label{lesson:TE-C403-pathakam}



\begin{quote}
Before the new one: say the Telugu for a representative, then the Telugu for coordination.
\end{quote}

\subsection*{You'll want to know: \te{పథకం}}
\textbf{\te{పథకం}} — \emph{pathakaṁ} — \textquotedblleft{}a scheme; a programme\textquotedblright{}.

\begin{quote}
\te{ప్రభుత్వం} \te{పేదల} \te{కోసం} \te{కొత్త} \te{పథకం} \te{తెచ్చింది}.
\end{quote}

\begin{quote}
\emph{prabhutvaṁ pēdala kōsaṁ kotta pathakaṁ teccindi.}
\end{quote}

\begin{quote}
The government brought in a new scheme for the poor.
\end{quote}

Plural: \te{పథకాలు}, \emph{pathakālu}.

Plural \te{పథకాలు}.

\begin{cousinweb}
\textbf{\te{పథకం}} (\emph{pathakaṁ}) has an origin that is not settled: it looks like a Sanskrit loan, but neither CDIAL nor Monier-Williams gives a source with this sense.
\end{cousinweb}

\begin{grammarlens}[title={What you've built}]
One more word for giving reasons and opinions.
\end{grammarlens}

\subsection*{Your turn}
\begin{itemize}
\raggedright
  \item \emph{Say it:} \emph{pathakaṁ}
  \item \emph{Say it:} \emph{pathakaṁ}, once more
  \item \emph{Say it:} \emph{prabhutvaṁ pēdala kōsaṁ kotta pathakaṁ teccindi.} — The government brought in a new scheme for the poor.
  \item \emph{Recall:} say the Telugu for a representative, then the Telugu for coordination, then say \emph{pathakaṁ} again
\end{itemize}

\begin{tcolorbox}[breakable,colback=teal!4,colframe=teal!35!black,title={Before you move on}]
What does \te{పథకం} mean? (\textquotedblleft{}A scheme; a programme\textquotedblright{}.) Say it once more.
\end{tcolorbox}
\section[\texorpdfstring{\te{సంక్షేమం} (saṅkṣēmaṁ)}{సంక్షేమం (saṅkṣēmaṁ)}]{\te{సంక్షేమం} (saṅkṣēmaṁ) — welfare}

\label{lesson:TE-C403-sanksemam}



\begin{quote}
Before the new one: say the Telugu for coordination, then the Telugu for a scheme.
\end{quote}

\subsection*{You'll want to know: \te{సంక్షేమం}}
\textbf{\te{సంక్షేమం}} — \emph{saṅkṣēmaṁ} — \textquotedblleft{}welfare\textquotedblright{}.

\begin{quote}
\te{పిల్లల} \te{సంక్షేమం} \te{కోసం} \te{ఆ} \te{సంస్థ} \te{పని} \te{చేస్తుంది}.
\end{quote}

\begin{quote}
\emph{pillala saṅkṣēmaṁ kōsaṁ ā saṁstha pani cēstundi.}
\end{quote}

\begin{quote}
That organisation works for the welfare of children.
\end{quote}

\te{సం}-, wholly, + \te{క్షేమం}, well-being, which you know.

\begin{cousinweb}
\textbf{\te{సంక్షేమం}} (\emph{saṅkṣēmaṁ}) is built from Sanskrit \emph{saṁ-}, \textquotedblleft{}together, wholly\textquotedblright{}, and \emph{kṣema}, \textquotedblleft{}ease, security, prosperity\textquotedblright{} (Monier-Williams), the word Telugu has as \te{క్షేమం}.
\end{cousinweb}

\begin{grammarlens}[title={What you've built}]
One more word for giving reasons and opinions.
\end{grammarlens}

\subsection*{Your turn}
\begin{itemize}
\raggedright
  \item \emph{Say it:} \emph{saṅkṣēmaṁ}
  \item \emph{Say it:} \emph{saṅkṣēmaṁ}, once more
  \item \emph{Say it:} \emph{pillala saṅkṣēmaṁ kōsaṁ ā saṁstha pani cēstundi.} — That organisation works for the welfare of children.
  \item \emph{Recall:} say the Telugu for coordination, then the Telugu for a scheme, then say \emph{saṅkṣēmaṁ} again
\end{itemize}

\begin{tcolorbox}[breakable,colback=teal!4,colframe=teal!35!black,title={Before you move on}]
What does \te{సంక్షేమం} mean? (\textquotedblleft{}Welfare\textquotedblright{}.) Say it once more.
\end{tcolorbox}
